\section{语料与数据}

\subsection{语料边界}

语料为《明朝那些事儿》全七部文本，共156章。项目从文本中整理1231个人物、563个地点、1106件事件与2324条关系。原始文本受版权约束，不随论文公开；统计底稿只保存汇总值、结构化条目与命中次数。这里的“明代思想史”始终带有语料限定：它是该书选择、复述和评论过的思想史片段。

\subsection{三层观察单位}

第一层是逐字语汇。我们调用统一接口对11个候选词做完全匹配，0命中词不进入论证。保留项中，“王守仁”命中538次，“心学”38次，“理学”29次，“书院”20次，“讲学”17次，“知行合一”12次，“致良知”2次。次数是字符串出现数，不是概念段落数。

第二层是关系。全图有12类关系，其中“师生同门”63条，占2.7\%。关系记录包含起点、终点、类别与章节；其中部分“其他”边由同事件共现自动推导。为避免共现把真实关系构成冲淡，正文以剔除推导边后的显式关系为主，同时保留全边口径作稳健性对照。

第三层是事件。我们只在事件名称和类型中匹配“心学、知行、格物、学派、论道、书院、讲学”，得到8件事件：6件被编码为人物生平，2件被编码为政治斗争。筛选规则窄而可复算，代价是不会捕获没有这些标记、但解释上可能相关的事件。

\begin{table}[!t]
\caption{本文的观察层与用途}
\label{tab:ideology-data}
\centering
\begin{tabular}{@{}lll@{}}
\toprule
层次 & 规模/口径 & 用途 \\
\midrule
原文语汇 & 11个非零候选词 & 确认叙事提及 \\
关系网络 & 2324条，12类 & 比较关系类别构成 \\
思想事件 & 8件，1488--1593 & 构造时间序列 \\
\bottomrule
\end{tabular}
\end{table}
